\section{The mean-cycle-length route fails in the reduced cubic class}
\label{sec:path-edge-half-mean-counterexample}

The mean-length condition remains stronger than the edge-half conclusion, even after the usual structural reductions.

\begin{proposition}\label{prop:path-edge-half-mean-counterexample}
There is a simple three-edge-connected cubic graph $G$ of girth five and order $66$ such that
\[
 c(G)=2\,160\,078\,826,\qquad
 \sum_{C\text{ a cycle}}|C|=107\,497\,549\,026.
\]
Consequently, its mean cycle length exceeds $3|V(G)|/4$.
The graph nevertheless satisfies the edge-half conclusion.
\end{proposition}

\begin{proof}
We first specify an eleven-vertex pole $N$.
Its vertices are $1,\ldots,11$, and its terminals are $4,7,10$.
List its edges in this order:
\[
\begin{split}
 &(1,5),(1,7),(1,8),(2,6),(2,8),(2,11),(3,7),(3,9),\\
 &(3,11),(4,8),(4,9),(5,10),(5,11),(6,9),(6,10).
\end{split}
\]
The terminals have degree two, and all other vertices have degree three.
The pole has rank five and girth five.
Its Hamiltonian cycle has vertex order
\[
 1,5,10,6,9,4,8,2,11,3,7.
\]
The remaining edges are $(1,8),(2,6),(3,9),(5,11)$.
In this cyclic order, the chord $(1,8)$ crosses each of the other three chords.

The pole is bridgeless because it has a Hamiltonian cycle.
Every two-edge cut consists of two Hamiltonian edges.
The intervening interval must contain both or neither endpoint of each chord.
Crossing chords cannot lie on different sides of such an interval.
Thus all eight chord endpoints lie on the same side.
The three terminals are pairwise nonconsecutive on the Hamiltonian cycle.
The other side must therefore be one terminal.
Every two-edge cut of $N$ isolates a terminal.

The cycle polynomial and the three terminal-path polynomials are
\[
\begin{aligned}
 C_N(z)&=5z^5+5z^6+4z^7+5z^8+5z^9+4z^{10}+z^{11},\\
 P_{47}(z)&=2z^3+z^5+7z^6+4z^7+2z^8+7z^9+3z^{10},\\
 P_{4,10}(z)=P_{7,10}(z)
 &=z^3+2z^4+2z^5+3z^6+5z^7+6z^8+5z^9+3z^{10}.
\end{aligned}
\]
For completeness, the following finite certificate verifies these counts without a graph search assumption.
Number the listed edges from one to fifteen.
A cycle-space basis consists of the five edge subsets
\[
 \{4,6,12,13,15\},\quad
 \{8,9,12,13,14,15\},\quad
 \{1,3,10,11,12,14,15\},\quad
 \{4,5,10,11,14\},\quad
 \{1,2,7,9,13\}.
\]
Their binary incidence vectors are independent, so they enumerate the complete cycle space by symmetric difference.
Number a basis sum by $s=\sum_{j=1}^5 2^{j-1}\varepsilon_j$.
For each parity class below, add the stated base edge set to every basis sum.
The table lists all nonzero states that are not connected.
All unlisted nonzero states are connected.
\[
\begin{array}{c|c|c|c}
\text{odd vertices}&\text{base edge set}&\text{disconnected indices }s&
 \text{their lengths}\hline
\varnothing&\varnothing&23,24&11,10\\
4,7&\{2,3,10\}&1,2,3,23,26,27&8,9,8,8,9,8\\
4,10&\{11,14,15\}&7,13,16,24,29&9,8,8,9,9\\
7,10&\{1,2,12\}&3,8,11,24,31&8,8,9,9,9
\end{array}
\]
This table follows directly by taking the indicated symmetric differences and checking their connected components.
Every edge coordinate is balanced in each parity class.
The total edge count over each class is therefore $15\cdot16=240$.
The table gives
\[
 c(N)=29,\quad \sum_C|C|=240-21=219,
 \qquad (p_{47},p_{4,10},p_{7,10})=(26,27,27),
\]
and the corresponding path-length sums are $(190,197,197)$.

Now replace every vertex of $K_{3,3}$ by a copy of $N$.
Order the three neighbours at each base vertex increasingly.
Attach their incident edges to ports $4,7,10$, in that order.
The resulting graph $G$ is simple and cubic, with $66$ vertices and $99$ edges.

Every cycle inside one pole has length at least five.
A cycle using edges between poles projects to a cycle of $K_{3,3}$.
Indeed, a cycle uses zero or two of the three boundary edges at each pole.
The projected cycle has length at least four.
Each internal terminal path has length at least three.
Every such crossing cycle therefore has length at least sixteen.
The internal five-cycles show that the girth of $G$ is exactly five.

The graph remains connected after deletion of any two edges.
If both deleted edges lie in one pole, every remaining pole component contains a terminal.
All these components connect through the other five poles, since $K_{3,3}$ minus one vertex is connected.
Otherwise each affected pole stays connected, because $N$ is bridgeless.
Contracting the poles then leaves $K_{3,3}$ with at most two edges removed, which is connected.
This proves three-edge connectivity.

Put
\[
 p_1=26,\quad p_2=p_3=27,\qquad
 \ell_1=190,\quad \ell_2=\ell_3=197,
\]
and set
\[
 A=p_1^2p_2^2p_3^2=359\,254\,116,\quad
 U=\sum_i p_i^2=2134,\quad V=\sum_i p_i\ell_i=15578.
\]
The nine base four-cycles contribute $U^2$ lifted cycles.
Their total length is $4U(U+V)$.
Each of the six base six-cycles uses each terminal-pair type twice.
They contribute $6A$ cycles and total length
\[
 36A+12A\sum_{i=1}^3\frac{\ell_i}{p_i}.
\]
There are also $6\cdot29$ cycles inside the poles, with total length $6\cdot219$.
Consequently
\[
\begin{aligned}
 c(G)&=6\cdot29+U^2+6A=2\,160\,078\,826,\\
 \sum_C|C|&=6\cdot219+4U(U+V)+36A
       +12A\sum_i\frac{\ell_i}{p_i}
       =107\,497\,549\,026.
\end{aligned}
\]
The exact difference is
\[
 4\sum_C|C|-3|V(G)|c(G)=2\,294\,588\,556>0.
\]

Finally, the edge $(2,11)$ lies in $15$ pole cycles and in $(12,13,13)$ paths of the three terminal types.
For the first pole, the three outside path coefficients are
\[
 (1405p_i+2A/p_i)_{i=1}^3
   =(27\,671\,462,26\,649\,351,26\,649\,351).
\]
The corresponding edge of $G$ therefore belongs to
\[
 15+12(27\,671\,462)+26(26\,649\,351)=1\,024\,940\,685
\]
cycles.
It is less than $c(G)/2=1\,080\,039\,413$.
Thus this graph obstructs the mean-length route but not the edge-half conclusion.
\end{proof}

\paragraph{Independent checks.}
The script \texttt{path-edge-half-mean-counterexample-check.py} uses only the Python standard library.
It enumerates the small pole by depth-first search and separately checks its affine-space certificate.
It reconstructs all $99$ edges and checks all $4\,851$ two-edge deletions.
It also computes every edge frequency from the exact gluing formulas.
Their sum equals the stated cycle-length sum.
A separate implementation by the random-host agent independently confirms these counts and structural checks.
